\documentclass{article}

\usepackage{PRIMEarxiv}

\usepackage[utf8]{inputenc} % allow utf-8 input
\usepackage[T1]{fontenc} % use 8-bit T1 fonts
\usepackage{hyperref} % hyperlinks
\usepackage{url} % simple URL typesetting
\usepackage{booktabs} % professional-quality tables
\usepackage{amsfonts} % blackboard math symbols
\usepackage{nicefrac} % compact symbols for 1/2, etc.
\usepackage{microtype} % microtypography
\usepackage{fancyhdr} % header
\usepackage{graphicx} % graphics
\usepackage{longtable} % tables that span pages
\usepackage{array} % column formatting helpers
\graphicspath{{media/}} % organize your images and other figures under media/ folder

%Header
\pagestyle{fancy}
\thispagestyle{empty}
\rhead{ \textit{ }}

% Update your Headers here
\fancyhead[LO]{Smartphone-Based Palm Oil Ripeness Detection}
% \fancyhead[RE]{Firstauthor and Secondauthor} % Firstauthor et al. if more than 2 - must use \documentclass[twoside]{article}

%% Title
\title
{Smartphone-Based Palm Oil Ripeness Detection Using Lightweight Vision Models
\thanks{\textit{\underline{Status}}: Final Year Project report (WIA3002), Faculty of Computer Science and
Information Technology, Universiti Malaya. \textbf{Supervisor}: Dr.\ Muhammad Shahreeza Safiruz Kassim.
\textbf{Co-Supervisor}: Ahmad Syazwan Ramli, Malaysian Palm Oil Board (MPOB).}
}

\author{
Nur Atiqah Maisarah Binti Mohamad Rasi \\
Faculty of Computer Science and Information Technology \\
Universiti Malaya \\
Kuala Lumpur, Malaysia \\
\texttt{atiqahmaisarah960@gmail.com} \\
}

\begin{document}
\maketitle

\begin{abstract}
Accurate ripeness grading of oil palm fresh fruit bunches (FFB) at the mill loading ramp is central to
oil extraction rate and to fair payment to suppliers, yet it is still performed by eye -- a judgement that
is subjective, fatigue-sensitive, and shifts with field lighting. This work builds a leakage-free,
four-source corpus of 18{,}167 labelled FFB images, split strictly by source photograph rather than by
crop, and trains three lightweight architectures -- MobileNetV2, EfficientNetB0, and the hybrid
convolution--transformer FastViT-T8 -- under one identical two-stage transfer-learning recipe, at three
seeds each. FastViT-T8 reaches \textbf{94.55\% $\pm$ 0.31\%} test accuracy on a held-out, source-balanced
split, clearing the project's $\geq$85\% target and fully repairing an unripe-class weakness (recall 0.38
$\rightarrow$ 0.940) inherited from a leaky, single-source FYP1 baseline. A leave-one-dataset-out (LODO)
evaluation, however, shows this accuracy does not transfer to an entirely unseen data source: accuracy
falls to \textbf{52.11\%}, a \textbf{42.44-point} gap that we attribute to source-specific labelling
conventions rather than to insufficient model capacity or training. The selected model converts losslessly
to a 6.1\,MB float16 TensorFlow Lite artifact using only built-in operators, and its accuracy varies by
just 1.02 points across simulated daylight, shadow, and low-light conditions. We further show that roughly
45\% of the apparent margin between FastViT-T8 and EfficientNetB0 is attributable to unequal fine-tuning
depth rather than to architecture, and that a fine-tuning-depth ablation run on an under-trained pipeline
produced a clean but false negative result -- a methodological caution of wider relevance to ablation
studies. On-device latency remains unverified pending a physical smartphone benchmark.
\end{abstract}

\keywords{oil palm fruit ripeness \and mobile deep learning \and TensorFlow Lite \and domain generalization \and lightweight vision transformers}

\section{Introduction}
\label{sec:intro}

\subsection{Background}

Oil palm (\textit{Elaeis guineensis}) is Malaysia's most economically significant agricultural commodity,
and the oil yield extracted from a fresh fruit bunch (FFB) depends heavily on when that bunch is harvested.
A bunch cut before maturity carries a low oil-to-bunch ratio; a bunch left too long accumulates free fatty
acids that degrade the oil's commercial grade. The window in which both quantity and quality are maximised
is narrow, and the decision to cut is made in the field, bunch by bunch, by a human grader.

That decision is currently made by eye. Graders judge ripeness from the colour of the fruitlets and the
number of loose fruits detached from the bunch, following criteria issued by the Malaysian Palm Oil Board
(MPOB) \citep{mpob2023grading}. Visual grading is fast and requires no equipment, but it is subjective:
judgement varies between graders, degrades with fatigue over a working day, and shifts with ambient
lighting -- the same bunch photographed in direct noon sun and in plantation shade presents noticeably
different colour. Inconsistent grading propagates directly into mill intake quality and into the price a
smallholder receives.

Automating this judgement is therefore a well-motivated computer-vision problem, and one that has
attracted substantial published attention \citep{suharjito2021mobile,herman2020visualattention}. The
constraint that shapes this project is not accuracy alone but \textit{where the computation happens}.
Malaysian oil palm smallholdings are frequently beyond reliable mobile data coverage, and the graders who
would use such a tool carry mid-range Android phones, not workstations. A system that requires a server
round-trip is unusable in the setting it is meant to serve. This project therefore targets a model that
runs entirely on-device, offline, within the memory and latency budget of a commodity smartphone.

\subsection{Problem statement}

Manual visual grading of oil palm FFB ripeness is subjective, inconsistent between graders, and sensitive
to field lighting conditions, which introduces avoidable variance into harvest timing and mill intake
quality. Existing automated approaches in the literature largely report accuracy on a single curated
dataset evaluated under laboratory or desktop conditions, and are seldom constrained by the memory,
latency, and offline-operation limits of the smartphone hardware a field grader actually carries.

This project's own preliminary work (FYP1) illustrates the gap concretely. A MobileNetV2
transfer-learning baseline trained on a public Roboflow dataset reached 68.14\% top-1 accuracy -- 17
points short of the project's $\geq$85\% acceptance bar. Diagnosis found the shortfall was not primarily a
modelling failure: 194 of the dataset's 205 unique source photographs had crops distributed across the
training, validation, \textit{and} test splits, so the reported accuracy was itself measured under
train/test leakage and overstated true generalization. The dataset contained only 205 unique source
photographs, all captured over two days, and the unripe class was represented by just 45 of them -- which
explains that model's 0.38 recall on that class. The problem is therefore twofold: build an evaluation
setting in which an accuracy number means what it claims to mean, and within that setting produce a model
that meets a smartphone deployment budget.

\subsection{Research objectives and questions}

Three research objectives were approved at FYP1 and are carried forward unchanged:

\begin{enumerate}
  \item \textbf{RO1} -- develop a smartphone-based oil palm FFB ripeness detection system using
  lightweight image classification models;
  \item \textbf{RO2} -- analyse and compare lightweight architectures -- MobileNet, EfficientNet-Lite, and
  Vision Transformers -- for ripeness classification;
  \item \textbf{RO3} -- evaluate the models under varying lighting conditions and optimize the
  best-performing model for smartphone deployment.
\end{enumerate}

Two substitutions were made within RO2 during execution, both forced by pretrained-weight availability
rather than by a research decision, and both documented in Section~\ref{sec:substitutions}:
EfficientNet-Lite0 was replaced by EfficientNetB0, and MobileViT by FastViT-T8. The objective's
requirement -- a lightweight CNN of each family plus a lightweight Vision Transformer variant -- is
preserved by both substitutions. These objectives translate into four research questions:

\begin{itemize}
  \item \textbf{RQ1} -- Can a lightweight vision model reach $\geq$85\% four-class FFB ripeness accuracy
  on a leakage-free, multi-source test split?
  \item \textbf{RQ2} -- Among a lightweight CNN (MobileNetV2), a scaled CNN (EfficientNetB0), and a
  hybrid convolution--transformer (FastViT-T8) trained under an identical recipe, which performs best,
  and is the margin larger than run-to-run variance?
  \item \textbf{RQ3} -- Does accuracy measured on held-out \textit{photographs} from sources present in
  training transfer to an entirely held-out \textit{source} -- a different camera, site, and labelling
  convention?
  \item \textbf{RQ4} -- Does the selected model survive conversion and quantization to TensorFlow Lite
  within the $\leq$100\,MB size and $\leq$2\,s latency budget, and does it remain stable across daylight,
  shadow, and low-light conditions?
\end{itemize}

Four hypotheses were formed against these questions. \textbf{H1}: replacing the leaky single-source FYP1
dataset with a leakage-free merged corpus from four independent sources raises achievable accuracy above
85\% and specifically repairs the unripe-class weakness, because that weakness is a data-diversity
artifact rather than a modelling limitation. \textbf{H2}: a hybrid convolution--transformer will
outperform the two convolutional baselines at comparable parameter cost, because global attention should
capture the gradual, whole-bunch colour transition between adjacent ripeness stages better than purely
local convolution. \textbf{H3}: accuracy on a held-out source will be materially lower than accuracy on
the main split, because the four source datasets encode different labelling conventions for the same
physical ripeness continuum. \textbf{H4}: the selected model meets the $\leq$100\,MB and $\leq$2\,s
budgets after quantization without a material accuracy penalty. H1 and H3 are supported by the results in
Section~\ref{sec:results}; H2 is supported with a substantial correction to its magnitude
(Section~\ref{sec:depth}); H4 is supported for model size and unresolved for on-device latency
(Section~\ref{sec:deployment}). One result was anticipated by none of these hypotheses and became the
project's central methodological finding: for most of this project's duration the binding constraint on
accuracy was neither the architecture nor the data but the training schedule
(Section~\ref{sec:schedule}).

\subsection{Scope, limitations, and contributions}

In scope: four ripeness classes (unripe, underripe, ripe, overripe); still images at $224\times224$
resolution; on-device inference with no server component; four public datasets as the training corpus.
Out of scope, stated here rather than left for the reader to discover: the empty-bunch and abnormal-fruit
categories present in one source and the ``decayed'' category present in another are excluded rather than
remapped, as they fall outside the four-class taxonomy; bunch detection and counting; and yield
estimation. Field images from MPOB via the co-supervisor were treated as an opportunistic bonus rather
than a dependency and did not arrive within the project timeline -- cross-source generalization is
therefore evidenced by the leave-one-dataset-out protocol of Section~\ref{sec:protocol} rather than by a
field trial. The Android application (Section~\ref{sec:appdesign}) was specified, designed, and
implemented; it runs the exported model end to end on an emulator and reproduces the Python pipeline's
probabilities to within the difference attributable to the platform's JPEG decoder
(Section~\ref{sec:ondevice}). What it has not had is a physical device: no latency figure, and no
photograph taken through the application in a plantation, exists. Field images could not be obtained
within the project period, which bounds what the deployment claims in this report can rest on and is
stated as such wherever they appear.

This project's contributions are: (1) a leakage-free merged corpus of 18{,}167 labelled FFB images from
four independently collected public sources, split by source photograph rather than by crop, with
per-(source, class) split balancing so that every source retains its own 70/15/15 proportions; (2) a
three-architecture comparison under a recipe held identical by construction, at three seeds each,
reporting run-to-run variance rather than single-run point estimates; (3) a quantified cross-source
generalization gap -- 42.44 points -- obtained by holding out each of the four sources entirely in turn,
together with an analysis identifying the failure mode as labelling-convention fitting rather than uniform
degradation; (4) a methodological finding of general applicability, that an ablation inherits the validity
of the recipe it is run under, demonstrated by a fine-tuning-depth control that produced a clean but false
negative result under an under-trained schedule and reversed when re-run correctly; and (5) a TensorFlow
Lite conversion study across three architectures and four quantization modes, establishing that
post-training full-integer quantization is unusable for the two more accurate architectures --
quantization robustness runs \textit{inversely} to the accuracy ranking.

\section{Related Work}
\label{sec:related}

\subsection{Oil palm ripeness grading: the task and its criteria}

MPOB's grading standard defines ripeness by fruitlet colour and by the count of loose fruits detached from
the bunch \citep{mpob2023grading}. The criteria are ordinal -- unripe, underripe, ripe, overripe form a
progression rather than a set of unrelated categories -- and the boundaries between adjacent stages are
conventions imposed on a continuous biological process rather than natural discontinuities. This matters
more than it first appears, and Section~\ref{sec:discussion} returns to it: two datasets can label the
same physical bunch differently without either being wrong, and a model trained across several such
datasets must reconcile conventions that do not agree.

\subsection{Deep learning for FFB ripeness classification}

Convolutional networks have displaced hand-designed colour and texture features for this task, with
reported accuracies commonly in the high eighties to mid nineties. Two published works bear directly on
this project's design decisions and are treated in detail.

\textbf{Suharjito et al.} -- the research group whose annotated bunch-pile dataset
\citep{suharjito2020bunchpiles} forms the largest component of this project's corpus -- published a
companion study applying ImageNet transfer learning to four lightweight CNNs for mobile FFB ripeness
classification \citep{suharjito2021mobile}. EfficientNetB0 was their best performer at \textbf{89.8\%
accuracy, 89.3\% after float16 TensorFlow Lite quantization, at 96\,ms per image}; MobileNetV1 was less
accurate at 81.1\% but fastest. Two specifics of their recipe are relevant here. First, they fine-tune at
\textbf{block-level granularity -- three unfrozen convolutional blocks} -- rather than by a fixed layer
count; this project adopts that scheme (Section~\ref{sec:recipe}), and Section~\ref{sec:depth} shows the
choice to be considerably more consequential than it appears. Second, they employ a fixed nine-angle
multi-crop augmentation scheme, which this project did \textit{not} adopt; it remains a candidate for
future work (Section~\ref{sec:futurework}).

\textbf{Harjanto et al.\ (2025)} \citep{harjanto2025yolov8seg} report up to \textbf{+6.58\% accuracy} from
segmenting the fruit bunch out of the background with a lightweight YOLOv8n-Seg model before classifying
with ShuffleNetV2. Their setting differs from this project's in ways that prevent the finding from being
assumed to transfer: a single source dataset, one object per image after manual curation, 2{,}400
hand-annotated segmentation masks, and no on-device latency or size constraint. This project scoped a
two-stage, cost-gated replication of the finding (Section~\ref{sec:futurework}) but did not execute it
within the project period.

\subsection{Lightweight architectures for mobile inference}

\textbf{MobileNetV2} \citep{sandler2018mobilenetv2} introduced inverted residual blocks with linear
bottlenecks and bounded ReLU6 activations, targeting mobile CPU inference directly. The bounded activation
range is a design decision made \textit{for} quantization, and Section~\ref{sec:deployment} observes its
consequence empirically: MobileNetV2 is the only one of the three architectures tested here that survives
post-training full-integer quantization intact.

\textbf{EfficientNet} \citep{tan2019efficientnet} derives its family by compound scaling of depth, width
and resolution from a mobile-sized base network. The EfficientNet-\textbf{Lite} variants were designed
specifically for edge deployment by removing squeeze-and-excitation blocks and replacing swish with ReLU6
-- both changes made for quantization compatibility. This project intended to use EfficientNet-Lite0 and
could not (Section~\ref{sec:substitutions}); the substitution to EfficientNetB0 retains the SE blocks and
SiLU activations that Lite removed, which Section~\ref{sec:deployment} shows is exactly why
EfficientNetB0's full-integer quantization degrades. The substitution is thus not merely a limitation but
an informative one.

\textbf{Vision Transformers on mobile.} ViTs \citep{dosovitskiy2021vit} model global relationships between
image patches through self-attention, at a computational cost and data appetite that are poor fits for
edge hardware. Mobile-oriented hybrids interleave convolution with attention to recover locality and
parameter efficiency: \textbf{MobileViT} \citep{mehta2022mobilevit} and \textbf{FastViT}
\citep{vasu2023fastvit} are the two considered here. MobileViT was named in this project's plan;
Section~\ref{sec:substitutions} documents why it could not be evaluated under the project's software stack
and why FastViT-T8 replaced it. The architectural argument motivating H2 is that the whole-bunch colour
gradient distinguishing adjacent ripeness stages is a global property of the image, which attention is
better positioned to capture than a purely local convolutional receptive field.

\subsection{Ordinal classification}

Because ripeness is ordinal, treating the four classes as unrelated categories discards structure:
confusing unripe with overripe is a worse error than confusing ripe with underripe, and categorical
cross-entropy is indifferent between them. Rank-consistent ordinal formulations such as \textbf{CORAL}
\citep{cao2020coral} decompose a $K$-class ordinal problem into $K{-}1$ rank-monotonic binary thresholds.
This project tested the approach directly; Section~\ref{sec:secondary} reports it as a negative result
with the mechanism identified.

\subsection{Datasets and the reproducibility problem}

Published FFB ripeness work is fragmented across small, independently collected datasets, and several
public datasets are re-exports or derivatives of one another. This project encountered the problem in a
concrete form. The ``new'' dataset sourced at the start of FYP2 proved to be the same public dataset
already used in FYP1, identified by identical per-class test counts. The archival copy of the bunch-piles
dataset turned out to be a Roboflow re-export whose class-id ordering differed from the ordering stated in
its own source publication, and in which 552 of roughly 4{,}160 unique frames appeared in more than one
split folder -- the same leakage defect that compromised FYP1, baked into the distribution. Both defects
are silent: neither raises an error, and both produce plausible-looking results.
Section~\ref{sec:construction} documents how each was detected and handled.

\textbf{Research gap.} Three gaps follow and define this project's contribution. First, accuracy figures
in this literature are predominantly single-source and single-run; the effect of a change of camera, site
and labelling convention is rarely quantified, and seed-to-seed variance is rarely reported at all.
Second, deployment claims are commonly made from framework-side latency measurements rather than from
converted, quantized models -- Section~\ref{sec:deployment} shows those to be a poor predictor. Third,
lightweight Vision Transformers are under-represented in FFB ripeness work relative to their prominence in
general mobile vision. This project addresses the first two directly and the third within the constraint
documented in Section~\ref{sec:substitutions}.

\section{Methodology}
\label{sec:methodology}

\subsection{Overall system architecture}

The system is organised as four layers, each layer's output forming the next layer's input. There is no
server component at any point: the trained model is embedded in the application and all inference happens
on the device -- Data Layer (four public sources merged and split by source photograph, with augmentation
applied on the fly to the training split only, plus a deterministic Lighting Stress-Test set) feeding the
Training Layer (one identical recipe across three architectures and three seeds, evaluated on accuracy,
F1, ROC-AUC, lighting robustness, and leave-one-dataset-out generalization) feeding the Export Layer
(TFLite conversion at four quantization settings, fidelity-checked against the Keras parent, then
device-benchmarked) feeding the App Layer (capture/upload UI, on-device TFLite inference, result display).

\subsection{App design}
\label{sec:appdesign}

\textit{Status: implemented. All four screens are built in Jetpack Compose and the flow runs end to end on
an Android~35 emulator; Section~\ref{sec:ondevice} reports what it does with the model.}
The interface follows a three-screen flow. \textbf{Home/Capture} offers two actions only -- \textit{Take
Photo} and \textit{Upload Image} -- with a status line confirming the model is loaded and the app works
offline. \textbf{Preview} shows the captured or selected image with \textit{Retake} and \textit{Detect
Ripeness}, so a poorly framed photo can be rejected before inference rather than silently misclassified.
\textbf{Result} shows the predicted class, its confidence, and, deliberately, a per-class confidence bar
for all four classes: because the underripe/ripe boundary is genuinely ambiguous
(Section~\ref{sec:discussion}), a bunch scoring 51\% ripe against 45\% underripe is materially different
information for a grader than one scoring 96\% ripe, and collapsing both to the single word ``Ripe'' would
hide exactly the uncertainty the grader needs to see. The application is designed to be offline by
construction (the \texttt{.tflite} model and \texttt{labels.txt} are bundled at build time and no network
permission is requested), to keep no account, upload, or history sync, to run inference on a background
thread with a determinate progress state, and to treat model-load and image-decode failures as defined
user-visible states rather than crashes.

Two aspects of the built application differ from the design above, both deliberately. The determinate
progress state was replaced by an indeterminate one: inference takes tens of milliseconds, so a percentage
would have been fiction, and the indicator is instead held for a 400\,ms floor so the step registers as a
step. More consequentially, the camera's framing brackets were specified as advisory decoration over a
preview whose \textit{whole} frame was to be classified; they are instead \textbf{load-bearing}, and the
capture is cropped to them before preprocessing. Section~\ref{sec:ondevice} gives the measurement that
forced that change. The offline property is enforced rather than asserted: the manifest requests no
network permission at all, so the model file and labels bundled at build time are the only possible
source of a prediction.

\subsection{Requirements}

Table~\ref{tab:nfr} lists the five non-functional requirements this project is evaluated against, together
with what was actually measured. In plain terms: the system has to be \textit{right} often enough to be
trusted, \textit{fast} enough not to interrupt fieldwork, \textit{small} enough to install on an ordinary
phone, \textit{steady} across the lighting a plantation actually offers, and free to run with no
connectivity.

\begin{table}[t]
\centering
\caption{Non-functional requirements and measured status.}
\label{tab:nfr}
\begin{tabular}{@{}llll@{}}
\toprule
Requirement & Target & Measured & Status \\
\midrule
Test-split accuracy & $\geq$85\% & \textbf{94.55\% $\pm$ 0.31\%} (FastViT-T8, $n{=}3$) & Met (scoped, \S\ref{sec:discussion}) \\
On-device latency & $\leq$2\,s & 47.2\,ms (converted, dev CPU, single-thread) & Not demonstrated \\
Deployed model size & $\leq$100\,MB & \textbf{6.1\,MB} (FastViT-T8, float16 TFLite) & Met, by $>$10$\times$ \\
Lighting stability & stable & 1.02\,$\pm$\,0.43\,pts across 4 conditions; 1.36\,pts deployed & Met \\
Offline / no per-use cost & required & No network permission requested; model bundled & Met \\
\bottomrule
\end{tabular}
\end{table}

NFR-1 and NFR-2 stand in tension with NFR-3 in principle, and the tension did not bind in practice: the
most accurate architecture also produced the smallest float16 artifact of the three (6.1\,MB against
EfficientNetB0's 7.7\,MB), because parameter count rather than accuracy drives file size. A genuine
trade-off does appear at the quantization level and is reported in Section~\ref{sec:deployment}: the
smallest and fastest setting is unusable for the selected architecture.

\subsection{Dataset construction}
\label{sec:construction}

\subsubsection{Sources}

Four independently collected public datasets were merged (Table~\ref{tab:sources}). The FYP1 Roboflow
dataset was retired as a training source and retained only as a documented negative result.

\begin{table}[t]
\centering
\caption{Source datasets merged into the training corpus.}
\label{tab:sources}
\begin{tabular}{@{}p{3.1cm}p{4.4cm}rp{4.4cm}@{}}
\toprule
Source & Origin & Rows & Notes \\
\midrule
Bunch Piles \citep{suharjito2020bunchpiles} & Mill-collected video frames, Indonesia; CC BY 4.0 & 12{,}583 & Ships as YOLO boxes; cropped to classification patches \\
Mendeley Ordinal \citep{kamal2024mendeley} & Central Kalimantan, industry-funded & 4{,}317 & 5 classes; ``Decayed'' dropped as out of scope \\
Kaggle \citep{hadeemadadum_kaggle} & Public supplement; CC BY 4.0 & 851 & Exact 4-class taxonomy match \\
Outdoor-Tenera Malaysia \citep{rosbi_outdoor_tenera} & Johor, N.\ Sembilan, Perak; outdoor lighting & 416 & No underripe class; damaged/empty dropped \\
\midrule
\textbf{Total} & & \textbf{18{,}167} & 9{,}473 unique source photographs \\
\bottomrule
\end{tabular}
\end{table}

The corpus is dominated by Bunch Piles (69\% of rows), and the Malaysia outdoor set -- the only
Malaysian-collected source, and therefore the closest proxy to the deployment setting -- contributes
2.3\%. Both facts bear on the interpretation of Section~\ref{sec:lodo}.

\subsubsection{Merge, taxonomy mapping, and split discipline}

All sources are mapped onto the locked four-class taxonomy; anything outside it is dropped rather than
remapped. Three defects were found in the source data during this process and are reported because each
was silent. First, \textbf{Bunch Piles' class-id order does not match its own publication} -- the Roboflow
re-export's \texttt{\_darknet.labels} gives a different 0--5 ordering than the paper states, so the
cropping script text-matches each label in Malay and English rather than hardcoding ids. Second,
\textbf{Bunch Piles frames appear in multiple split folders}: 552 of $\sim$4{,}160 unique frames occur in
more than one of the export's own train/valid/test directories, so every source's own split is ignored and
the split is rebuilt from scratch across the merged corpus. Third, \textbf{byte-identical duplicate images
exist within sources}, including two pairs crossing the Malaysia set's own train/test boundary; every
image is MD5-hashed and exact-content duplicates dropped (11 from Malaysia, 6 pairs from Kaggle). All
whole-photo sources were also cross-checked against each other by MD5: \textbf{zero images are shared
across datasets}.

Splitting groups every row by (source dataset, source photo) and assigns whole groups to train/valid/test
at 70/15/15, so every crop from one bunch-pile frame stays in one split. \textbf{Zero source-photo leakage
is verified programmatically} -- every group's rows land in exactly one split (Table~\ref{tab:split}).

\begin{table}[t]
\centering
\caption{Class counts by split after the corrected, per-(source, class) balanced split.}
\label{tab:split}
\begin{tabular}{@{}lrrrrr@{}}
\toprule
Split & Overripe & Ripe & Underripe & Unripe & Total \\
\midrule
Train & 2{,}580 & 3{,}648 & 3{,}259 & 3{,}229 & 12{,}716 \\
Valid & 553 & 784 & 699 & 693 & 2{,}729 \\
Test  & 552 & 780 & 698 & 692 & 2{,}722 \\
\bottomrule
\end{tabular}
\end{table}

\textbf{A split-balancing bug was found and fixed.} The first version of the split assigned groups against
a global per-class target. Because Bunch Piles' groups are large (many crops per frame) while every other
source contributes single-image groups, and the greedy assignment sorts by group size descending and fills
whichever split is furthest below target, Bunch Piles was funnelled into train -- 82--94\% of every
class's train split against as little as 14\% of valid/test. Models were being evaluated on a different
source mix than they trained on, for every class. The fix computes targets per (source dataset, class), so
each source retains its own 70/15/15 proportions regardless of group size. The correction was worth
roughly ten points of measured accuracy on the baseline model (63.71\% $\rightarrow$ 73.73\%), and every
result in this report is from the corrected split.

\subsubsection{Architecture substitutions}
\label{sec:substitutions}

Both substitutions below were forced by pretrained-weight availability under the project's software stack
(TensorFlow 2.20 / Keras 3.15), not chosen for research reasons.

\textbf{EfficientNet-Lite0 $\rightarrow$ EfficientNetB0.} The only available pretrained Lite0 weights are
published in the legacy TF1-Hub format, which raises a \texttt{ValueError} on \texttt{trainable=True} and
therefore cannot be fine-tuned. Three integration routes were tested and all failed for reasons independent
of the fine-tuning limitation. Since every other architecture in the comparison uses two-stage
frozen-then-fine-tuned training, a frozen-only Lite0 would confound architecture with training
methodology. EfficientNetB0 is Keras-3 native, supports the identical recipe, and belongs to the same
lightweight CNN family.

\textbf{MobileViT $\rightarrow$ FastViT-T8.} Five independent integration routes for MobileViT each failed
for a distinct, reproducible reason, the most fundamental being a \texttt{KerasTensor} type error inside
\texttt{transformer\_pre\_process} at the patch-folding step -- an operation constitutive of the MobileViT
architecture, so not workable around without reimplementing the model. To establish whether the obstacle
was specific to MobileViT, fifteen lightweight transformer variants were swept under identical conditions:
\textbf{only FastViT built successfully}; MobileViT V2, LeViT, EfficientViT, EdgeNeXt and TinyViT all
failed identically. The limitation is a property of the current Keras 3 ecosystem for mobile transformers,
not of this project's configuration. FastViT-T8 was then validated against six criteria before any
training time was committed -- including lossless save/reload round-trips, the criterion on which an
earlier candidate route had failed. At 3.20M parameters it sits between MobileNetV2 (2.26M) and
EfficientNetB0 (4.05M), so the comparison is not confounded by model scale.

\subsubsection{Training recipe}
\label{sec:recipe}

All three architectures are driven through one shared two-stage schedule (Table~\ref{tab:recipe}).
Fine-tuning depth is specified in \textbf{blocks rather than layers}, following
\citet{suharjito2021mobile}, because a fixed layer count represents different proportional depths in
different architectures. Section~\ref{sec:depth} shows this choice to be far more consequential than a
routine hyperparameter. Label smoothing at 0.15 is retained from FYP1 for a task-specific reason: the
ripe/underripe boundary is genuinely subjective, and label smoothing makes the model appropriately less
confident on borderline cases rather than forcing a hard call.

\begin{table}[t]
\centering
\caption{Shared training recipe (v2) across all three architectures.}
\label{tab:recipe}
\begin{tabular}{@{}p{3.6cm}p{9.8cm}@{}}
\toprule
Element & Setting \\
\midrule
Input & $224\times224\times3$, batch size 32 \\
Head & GlobalAveragePooling $\rightarrow$ Dropout(0.2) $\rightarrow$ Dense(4, softmax) \\
Loss & Categorical cross-entropy, label smoothing 0.15 \\
Class imbalance & Balanced class weights applied as per-sample weights \\
Stage 1 (frozen) & Backbone frozen, head only, lr $10^{-3}$, cap 30--50 epochs \\
Stage 2 (fine-tune) & Last $K{=}3$ blocks unfrozen, lr $10^{-4}$ under cosine decay to $\sim$0, cap 40 epochs \\
Early stopping & Patience 8 on validation loss, best weights restored \\
BatchNorm & Pinned to inference mode throughout, including Stage 2 \\
Augmentation & Train split only: horizontal flip, $\pm$10\% rotation/zoom/translation, $\pm$20\% brightness and contrast, applied on-the-fly in $[0,255]$ space before rescaling \\
Seeds & 42, 43, 44 for every reported configuration \\
\bottomrule
\end{tabular}
\end{table}

\subsection{Evaluation protocol}
\label{sec:protocol}

Two splits are used and must not be confused (Table~\ref{tab:protocol}). Metrics reported: overall
accuracy, per-class precision/recall/F1, macro F1, one-vs-rest ROC-AUC, confusion matrices, parameter
count, artifact size, and inference latency. For LODO folds, macro F1 is computed over the classes
actually present in the held-out source -- the Malaysia set contains no underripe images, and averaging
over the full label set would penalise that fold for a class it could not have been tested on.

\begin{table}[t]
\centering
\caption{The two evaluation splits used in this project.}
\label{tab:protocol}
\begin{tabular}{@{}p{2.6cm}p{5.2cm}p{5.6cm}@{}}
\toprule
 & Main split & Leave-one-dataset-out (LODO) \\
\midrule
Construction & All four sources pooled, split by source photograph; every source appears in train, valid and test & Each source held out entirely in turn; the other three form the training pool (four runs) \\
Question & How accurate is the model on an unseen photograph from these sources? & Does the model work on a camera, site, and labelling convention it has never encountered? \\
Role & Headline figure against NFR-1; selects the winning architecture & Primary generalization evidence; the honest estimate of field performance \\
Cost & All three architectures $\times$ 3 seeds & Winner only: 4 folds $\times$ 3 seeds \\
\bottomrule
\end{tabular}
\end{table}

\subsection{Lighting robustness protocol}

The lighting stress-test set is generated once from the clean test split and versioned as its own
artifact, never randomly regenerated. Three conditions are applied to every one of the 2{,}722 test
images: \textbf{darkened} (brightness $\times$0.7, simulating low light), \textbf{brightened} (brightness
$\times$1.3, simulating harsh daylight and overexposure), and \textbf{shadow} (a soft-edged elliptical
patch at 50\% brightness, its position and size seeded deterministically from a hash of the image path, so
the effect varies image-to-image like a real cast shadow but reproduces exactly on every run). This adds a
shadow condition over FYP1's brightness-only test so that all three lighting conditions named in NFR-4 are
actually covered; the resulting figures are not numerically comparable with FYP1's lighting numbers and
are not presented as such.

\subsection{Reproducibility measures}

Every run writes a configuration file (including a \texttt{recipe\_version} field), per-epoch history,
metrics, a classification report, a confusion matrix, and lighting stress-test metrics to its own results
directory. The aggregation scripts refuse to pool runs across recipe versions or across compute backends,
on the grounds that a v1 run and a v2 run are not replicates of the same process -- nor are a CPU run and
a GPU run. All reported configurations use three seeds, and every headline figure is quoted as mean $\pm$
sample standard deviation rather than as a single run.

\section{Results and Analysis}
\label{sec:results}

All results in this section were produced under training recipe \textbf{v2} on a single NVIDIA RTX 4050
Laptop GPU (TensorFlow 2.20 / Keras 3.15 / Python 3.12). Where an earlier v1 figure is quoted it is
explicitly labelled as superseded; Section~\ref{sec:schedule} explains why the distinction is
load-bearing.

\subsection{Headline result}

\textbf{FastViT-T8, fine-tuned at three blocks, reaches 94.55\% $\pm$ 0.31\% test accuracy on the held-out
test split ($n{=}3$ seeds), with macro F1 0.9467 $\pm$ 0.0035 and one-vs-rest ROC-AUC 0.9916 $\pm$
0.0002.}

\begin{quote}
\textit{Required scoping -- reproduced wherever the headline appears.} The $\geq$85\% accuracy requirement
is met on the held-out test split (94.55\% $\pm$ 0.31\%). It is \textbf{not} claimed to hold under source
shift: leave-one-dataset-out accuracy is 52.11\%, a 42.44-point generalization gap, and that figure is the
honest estimate of performance on a mill this model has never seen.
\end{quote}

All three architectures clear the 85\% bar on the main split, so \textbf{RQ1 is answered affirmatively and
H1 is supported}. The FYP1 unripe-class failure is fully repaired: recall on unripe rises from 0.38 to
0.940 $\pm$ 0.005, confirming that the weakness was a data-diversity artifact rather than a modelling
limitation.

\subsection{Architecture comparison}
\label{sec:archcomparison}

\begin{table}[t]
\centering
\caption{Three-architecture comparison, recipe v2, mean $\pm$ std over 3 seeds. Source:
\texttt{multiseed\_v2\_summary.md}.}
\label{tab:archcomparison}
\begin{tabular}{@{}lrrr@{}}
\toprule
Metric & MobileNetV2 & EfficientNetB0 & \textbf{FastViT-T8} \\
\midrule
Test accuracy (\%) & 89.62 $\pm$ 0.15 & 91.26 $\pm$ 0.39 & \textbf{94.55 $\pm$ 0.31} \\
ROC-AUC (OvR) & 0.9800 $\pm$ 0.0018 & 0.9861 $\pm$ 0.0012 & \textbf{0.9916 $\pm$ 0.0002} \\
Macro F1 & 0.8983 $\pm$ 0.0015 & 0.9137 $\pm$ 0.0042 & \textbf{0.9467 $\pm$ 0.0035} \\
Recall -- unripe & 0.9397 $\pm$ 0.0107 & 0.9427 $\pm$ 0.0021 & 0.9400 $\pm$ 0.0053 \\
Recall -- underripe & 0.7933 $\pm$ 0.0150 & 0.7960 $\pm$ 0.0108 & \textbf{0.9033 $\pm$ 0.0161} \\
Recall -- ripe & 0.9207 $\pm$ 0.0055 & 0.9553 $\pm$ 0.0068 & \textbf{0.9663 $\pm$ 0.0021} \\
Recall -- overripe & 0.9363 $\pm$ 0.0095 & 0.9623 $\pm$ 0.0064 & \textbf{0.9763 $\pm$ 0.0086} \\
Lighting spread (pts) & 5.19 $\pm$ 0.88 & 2.45 $\pm$ 0.37 & \textbf{1.02 $\pm$ 0.43} \\
Parameters (M) & 2.26 & 4.05 & 3.20 \\
Size, pre-quantization (MB) & 20.9 & 33.9 & 30.2 \\
\bottomrule
\end{tabular}
\end{table}

\begin{figure}[t]
\centering
\includegraphics[width=0.85\linewidth]{09_multiseed_architecture_comparison.png}
\caption{Three-architecture comparison with seed-to-seed spread (recipe v2).}
\label{fig:archcomparison}
\end{figure}

FastViT-T8 leads on every accuracy metric by margins several times the seed-to-seed spread, so \textbf{RQ2
is answered: the ranking is real, not noise}. The margin over EfficientNetB0 at the same nominal $K$ is
3.29 points -- but Section~\ref{sec:depth} shows that figure to overstate the architectural contribution
substantially, and it must not be quoted without that correction. The largest single per-class gain is on
\textbf{underripe}, the class that had been this project's persistent weak point: 0.79 for both CNNs
against 0.903 for FastViT. That the hybrid architecture's advantage concentrates on the hardest ordinal
boundary is consistent with the mechanism proposed in H2 -- global attention capturing a whole-bunch
colour gradient -- though this report does not claim to have isolated that mechanism.

The v2 campaign also recorded framework-side inference times of 141.9 / 161.3 / 156.6\,ms for
MobileNetV2 / EfficientNetB0 / FastViT-T8. Section~\ref{sec:deployment} establishes those figures to be a
poor predictor of deployed latency, and they are reported here only to be superseded -- the converted
models measure 12.1 / 22.4 / 47.7\,ms, a 3.9$\times$ spread where the framework figures showed
1.13$\times$, and in a different rank order.

\subsubsection{The training-schedule correction: the project's central methodological finding}
\label{sec:schedule}

For most of this project's duration every reported accuracy was an artifact of an under-trained pipeline,
and the defect was invisible in the results. The original schedule (recipe v1) capped training at 10
frozen plus 10 fine-tuning epochs with a constant $10^{-5}$ fine-tuning learning rate. Reading back the
saved per-epoch history of all fifteen GPU runs afterwards showed fine-tuning (Stage 2) early-stopping in
\textbf{0 of 15} runs, against 6 of 15 for the frozen Stage 1; validation loss was still falling at the
final epoch in 13 of 15 runs, and the best validation accuracy occurred at the final epoch in 11 of 15.
Fine-tuning hit its epoch cap in \textbf{every single run}. That early stopping \textit{did} fire in Stage
1 for six of those runs is what makes the diagnosis conclusive rather than a suspected dead callback: the
frozen head genuinely plateaus and the callback correctly ends it, while the fine-tuning stage never got
the chance.

Recipe v2 raises the caps well clear of where runs actually stop (30--50 frozen, 40 fine-tuning), raises
early-stopping patience from 4 to 8, and raises the fine-tuning learning rate to $10^{-4}$ under cosine
decay. On \textbf{unchanged data and unchanged splits}, a single FastViT seed moved from 78.69\% to
93.98\%, since confirmed across all three architectures and three seeds each (Table~\ref{tab:v1v2}).

\begin{table}[t]
\centering
\caption{Recipe v1 vs.\ v2 test accuracy, by architecture.}
\label{tab:v1v2}
\begin{tabular}{@{}lrrr@{}}
\toprule
Architecture & Recipe v1 (\%) & Recipe v2 (\%) & $\Delta$ \\
\midrule
MobileNetV2 & 73.61 $\pm$ 0.59 & 89.62 $\pm$ 0.15 & +16.01 \\
EfficientNetB0 & 76.45 $\pm$ 0.90 & 91.26 $\pm$ 0.39 & +14.81 \\
FastViT-T8 & $\sim$78.7 (single seed) & 94.55 $\pm$ 0.31 & $\approx$+15.9 \\
\bottomrule
\end{tabular}
\end{table}

Two variables move together in v2 -- training budget and learning rate -- so this is an upgraded recipe
applied uniformly, not a controlled single-variable ablation. The \textit{architecture} comparison remains
fair because every architecture receives the same v2. The finding that matters beyond this project is in
Section~\ref{sec:depth}: an ablation run under v1 inherited v1's invalidity, and produced a clean result
that was false.

\subsection{Fine-tuning depth: how much of the architecture margin is really architecture?}
\label{sec:depth}

FastViT-T8 exposes ten blocks; MobileNetV2 and EfficientNetB0 expose sixteen each. $K{=}3$ therefore
unfreezes \textbf{28.0\%} of FastViT's layers but only \textbf{19.4\%} of EfficientNetB0's and 18.8\% of
MobileNetV2's. The same nominal setting is not the same proportional depth, so the 3.29-point margin may
reflect fine-tuning depth rather than architecture. A control was run to separate them: EfficientNetB0 at
$K{=}4$ -- 25.7\% of its layers, the closest available match to FastViT's 28.0\% -- at three seeds,
everything else identical (Table~\ref{tab:depth}).

\begin{table}[t]
\centering
\caption{Fine-tuning depth control. Source: \texttt{effnet\_depth\_k4\_v2\_summary.md}.}
\label{tab:depth}
\begin{tabular}{@{}lrrl@{}}
\toprule
Configuration & Layers unfrozen & Test accuracy (\%) & Per-seed \\
\midrule
EfficientNetB0, $K{=}3$ & 19.4\% & 91.26 $\pm$ 0.39 & 91.55 / 91.40 / 90.82 \\
\textbf{EfficientNetB0, $K{=}4$} & \textbf{25.7\%} & \textbf{92.73 $\pm$ 0.37} & 92.73 / 92.36 / 93.09 \\
FastViT-T8, $K{=}3$ (selected) & 28.0\% & 94.55 $\pm$ 0.31 & -- \\
\bottomrule
\end{tabular}
\end{table}

\textbf{At comparable fine-tuning depth, FastViT-T8's advantage over EfficientNetB0 is 1.82 points rather
than the 3.29 points implied by the same-$K$ comparison. Approximately 45\% of the apparent architectural
advantage is attributable to fine-tuning depth rather than to architecture.} The residual 1.82 points
exceeds the combined run-to-run spread of $\pm$0.68, so the architecture conclusion survives -- H2 is
supported, at roughly half the magnitude it first appeared to have. The depth effect is not confined to
aggregate accuracy: at $K{=}4$, EfficientNetB0's underripe recall rises 0.796 $\rightarrow$ 0.831 and its
lighting spread tightens from 2.45 to 1.33 points. For scale, the accuracy gap between EfficientNetB0 and
MobileNetV2 -- two architectures with identical block counts, directly comparable at a shared $K$ -- is
1.64 points. \textbf{A single step of fine-tuning depth is worth about as much as the difference between
two architectures.}

\textbf{The falsified control: a methodological result, not an admission.} This same control was run once
before, under recipe v1, and produced the opposite result: EfficientNetB0 at $K{=}4$ scored 0.42 points
\textit{below} $K{=}3$, comfortably within run-to-run variation. That measurement supported a conclusion --
that depth sensitivity was specific to the hybrid transformer and the architecture margin survived the
control undiminished -- which is now known to be an artifact. The v1 schedule halted every run before
convergence, so the additional unfrozen layers had no remaining budget in which to adapt. \textbf{The
control could not have detected a depth effect had one been present.} The episode generalises beyond this
project: an ablation inherits the validity of the recipe it is run under. A control executed on an
under-trained pipeline measures the under-training, not the variable it was designed to isolate. Because
the defect \textit{suppressed} the treatment effect rather than producing an implausible number, the
result looked clean and was accepted; it was caught only because the recipe was later corrected for
unrelated reasons and every dependent conclusion was re-derived rather than carried forward.

\textbf{Residual limitation.} Depth is still not exactly matched and cannot be: EfficientNetB0's block grid
steps 19.4\% $\rightarrow$ 25.7\% $\rightarrow$ 31.6\% while FastViT's steps 15.4\% $\rightarrow$ 28.0\%,
leaving a residual mismatch of about 2.3 points with the CNN on the shallower side. That mismatch is not
neutral -- EfficientNetB0 was still gaining as depth increased rather than having saturated -- so
\textbf{the 1.82-point margin should be read as an upper bound on the architectural contribution, not a
point estimate.} A $K{=}5$ run (31.6\%), bracketing FastViT from above, is the natural next control and
was not carried out.

\subsection{Cross-source generalization (LODO)}
\label{sec:lodo}

Each source was held out entirely in turn, FastViT-T8 trained on the remaining three, and evaluated purely
on the held-out source: four folds $\times$ three seeds (Table~\ref{tab:lodo}, Figure~\ref{fig:lodo}).

\begin{table}[t]
\centering
\caption{Leave-one-dataset-out (LODO) generalization, FastViT-T8, recipe v2. Source: \texttt{lodo\_v2\_summary.md}.}
\label{tab:lodo}
\begin{tabular}{@{}lrrrr@{}}
\toprule
Held-out source & Classes & Test accuracy (\%) & ROC-AUC & Macro F1 (present classes) \\
\midrule
Bunch Piles & 4 & 40.98 $\pm$ 2.33 & 0.7153 $\pm$ 0.0205 & 0.4128 $\pm$ 0.0194 \\
Kaggle & 4 & 50.80 $\pm$ 0.78 & 0.8295 $\pm$ 0.0023 & 0.4972 $\pm$ 0.0051 \\
Mendeley Ordinal & 4 & 55.75 $\pm$ 0.40 & 0.7900 $\pm$ 0.0097 & 0.5474 $\pm$ 0.0073 \\
Malaysia Outdoor & 3 & 60.90 $\pm$ 0.37 & n/a & 0.6503 $\pm$ 0.0125 \\
\midrule
\textbf{Mean across folds} & & \textbf{52.11} & & \\
\bottomrule
\end{tabular}
\end{table}

\begin{figure}[t]
\centering
\includegraphics[width=0.85\linewidth]{11_lodo_generalization_gap.png}
\caption{Main-split accuracy vs.\ leave-one-dataset-out accuracy, per fold.}
\label{fig:lodo}
\end{figure}

\textbf{The generalization gap is 42.44 points} (94.55\% main split against 52.11\% LODO mean).
\textbf{RQ3 is answered and H3 is supported, at a magnitude larger than anticipated.} Three properties of
this result require statement rather than summary.

\textbf{It is not another under-training artifact.} Under v1, 0 of 15 runs early-stopped; under v2,
\textbf{10 of 12 fine-tuning stages early-stopped.} The low cross-source number cannot be explained away
by training budget the way the v1 main-split numbers could. The honest reading is that a properly trained
model fits each training source's labelling convention \textit{harder} rather than generalizing past it --
correcting the recipe bought roughly 18 points on held-out photographs and gave back roughly 4 points on
an entirely held-out source (the v1 LODO mean was 55.96\%, against a 76.45\% main split: a 20.49-point gap
that has now widened to 42.44).

\textbf{The failure mode is convention-fitting, not uniform degradation.} Under v2 the model concentrates
its cross-source predictions on ripe and overripe. On held-out Kaggle, ripe recall rose 0.659
$\rightarrow$ 0.921 and overripe 0.785 $\rightarrow$ 0.912 between v1 and v2, while \textbf{underripe
collapsed from 0.406 to 0.067}. On held-out Mendeley the same pattern appears (ripe 0.621 $\rightarrow$
0.817, underripe 0.128 $\rightarrow$ 0.063). Held-out Bunch Piles moves the other way (underripe 0.199
$\rightarrow$ 0.547, unripe 0.418 $\rightarrow$ 0.140), consistent with the same mechanism rather than
against it: \textbf{each fold snaps to whichever convention the remaining three sources agree on.}

\textbf{Two claims about underripe must be kept apart.} Main-split underripe weakness was under-training,
and v2 fixed it ($\approx$0.50 $\rightarrow$ 0.903). Cross-source underripe weakness is a
labelling-convention mismatch between sources, and v2 made it \textit{worse}. These are backed by different
campaigns pointing in opposite directions and must not be conflated.

\textbf{Confounds to state.} Fold training-set sizes differ substantially -- holding out Bunch Piles leaves
$\sim$4.7k training images against $\sim$15.1k for Malaysia Outdoor -- so the weak Bunch Piles fold
confounds domain shift with reduced training data; the other three folds isolate domain shift more
cleanly. The v1 LODO campaign ran EfficientNetB0 while the v2 campaign runs FastViT-T8, so recipe and
architecture changed together and the $-$3.85 delta is not attributable to the recipe alone. What is safe
to claim under any attribution is the \textbf{direction}: cross-source accuracy did not improve.

\subsection{Mobile deployment: conversion, quantization, and size}
\label{sec:deployment}

All three v2 architectures at seed 42 were converted to TensorFlow Lite at four quantization settings,
with each converted model's predictions compared against its unquantized Keras parent on a 200-image
sample (Table~\ref{tab:tflite}).

\begin{table}[t]
\centering
\caption{TensorFlow Lite conversion and quantization results (seed 42). Latency is single-threaded x86
development CPU. Source: \texttt{tflite\_export\_summary.md}.}
\label{tab:tflite}
\begin{tabular}{@{}llrrrr@{}}
\toprule
Model & Mode & Size (MB) & Accuracy ($n{=}200$) & Agreement w/ parent & Latency (ms) \\
\midrule
FastViT-T8 & float32 & 12.1 & 0.920 & 0.990 & 47.7 \\
\textbf{FastViT-T8} & \textbf{float16} & \textbf{6.1} & \textbf{0.925} & 0.970 & \textbf{47.2} \\
FastViT-T8 & dynamic int8 & 3.3 & 0.930 & 0.970 & 114.3 \\
FastViT-T8 & full int8 & 3.5 & \textbf{0.260} & 0.265 & 136.7 \\
EfficientNetB0 & float16 & 7.7 & 0.910 & 1.000 & 23.6 \\
EfficientNetB0 & full int8 & 4.7 & \textbf{0.645} & 0.660 & 20.3 \\
MobileNetV2 & float16 & 4.3 & 0.905 & 1.000 & 12.0 \\
MobileNetV2 & full int8 & 2.6 & 0.915 & 0.960 & 13.4 \\
\bottomrule
\end{tabular}
\end{table}

\textbf{All twelve conversions succeeded using built-in TensorFlow Lite operators alone -- none required
the Flex delegate.} This matters because FastViT-T8 is a hybrid convolution--transformer assembled from
custom layers, and a Flex-dependent conversion would have obliged the application to embed a slice of the
full TensorFlow runtime, costing tens of megabytes. The selected architecture is deployable against a
stock TFLite interpreter, and the deployability risk is closed.

\textbf{Recommended artifact: FastViT-T8 at float16 -- 6.1\,MB, 47.2\,ms, at a measured accuracy cost of
0.44 points.} Scored on the full 2,722-image test split rather than the 200-image fidelity sample, the
float16 artifact reaches 94.12\% against its Keras parent's 94.56\% at the same seed
(Section~\ref{sec:lightingresults}); the sample's apparent one-point \textit{gain} was sampling noise,
and the deployed model is slightly worse than the trained one. Half the
float32 size for no latency penalty. Size is nowhere near binding (6.1\,MB against 100\,MB), so there is
no reason to accept dynamic int8's 2.4$\times$ latency penalty to save a further 2.8\,MB. \textbf{NFR-3 is
met with a margin of more than an order of magnitude.}

\textbf{Post-training full-integer quantization is unusable for the two more accurate architectures, and
the failure tracks architecture rather than size.} FastViT collapses to 0.260 -- chance on four classes --
and EfficientNetB0 to 0.645, while MobileNetV2 is untouched at 0.915. This is the textbook pattern:
FastViT's normalization and attention activations and EfficientNet's SiLU plus squeeze-and-excitation
branches both produce dynamic ranges that a single per-tensor integer scale cannot represent, whereas
MobileNetV2's bounded ReLU6 design was built to quantize. \textbf{Quantization robustness runs inversely
to the accuracy ranking: the least accurate architecture is the most deployable.} All three fully-quantized
models converted without error and produced the smallest files of any setting tested -- a benchmark
measuring only conversion success, file size and latency would have selected the fully-quantized model for
all three architectures, including the one performing at chance; the failure is detectable only by
comparing the converted model's predictions against those of the model it was derived from.

\textbf{What this does not establish.} Seed 42 only, not the three-seed protocol. The accuracy column is
measured on a 200-image sample sized to compare converted against unconverted predictions -- on that
sample the three architectures sit within one point of each other while the full test set separates them
by five, so it must not be read as an architecture comparison. Latency is single-threaded x86; integer
arithmetic frequently outperforms float16 on ARM via XNNPACK or NNAPI, so the float16 recommendation could
reverse on a phone. \textbf{NFR-2 remains undemonstrated: a physical Android benchmark was not carried
out.}

\subsection{On-device implementation and input fidelity}
\label{sec:ondevice}

The application described in Section~\ref{sec:appdesign} was built and runs the float16 artifact end to
end. This section reports what that exercise established about the model's \textit{input}, which turned
out to matter more than anything it established about its output: the largest single accuracy effect
measured anywhere in this project was found here, in the camera's framing, and not in the model.

\textbf{The application reproduces the Python pipeline, and the residual difference is the JPEG decoder.}
Three measurements separate the possible causes rather than asserting a conclusion. First, fed the exact
float32 tensors the Python pipeline was fed, on twenty stratified test-split fixtures bundled with the
application, the on-device interpreter reproduces the reference probabilities to within
\textbf{$1\times10^{-4}$} -- establishing that interpreter wiring, buffer layout, output shape and label
order are all correct. Second, a JVM unit test confirms the Kotlin bilinear resize matches
\texttt{tf.image.resize} under half-pixel centres without antialiasing to better than $1\times10^{-3}$ at
800 sampled positions, removing the resize as a candidate. Third, starting instead from the same JPEG
files and running the application's own decode--resize--scale path, probabilities move by up to
\textbf{0.0424}, with the \textbf{predicted class unchanged on all twenty}. With the first two fixed, the
third is attributable to image decoding alone: Android cannot call \texttt{tf.io.decode\_jpeg}, and
decoding the same fixtures on the desktop with Pillow -- an independent libjpeg configuration --
disagrees with TensorFlow's decoder by up to 16 levels of 255 per pixel (mean 1.8) and reproduces the
\textit{same} worst-case probability shift, on the same image, to seven decimal places (0.042418). The
difference is the decoder, not the application. The bundling script now performs this comparison as part
of building the app bundle and derives the instrumented test's tolerance from it (0.064) rather than from
an assumed value.

\textbf{Over the full test split the decoder costs nothing measurable.} The twenty fixtures bound
the decoder's effect image by image; its effect on accuracy over the whole split can be read off two
runs already in this report. The same float16 artifact (identical SHA-256) scores \textbf{94.16\%} on
all 2,722 test images decoded with Pillow -- the baseline of the framing study below -- and
\textbf{94.12\%} decoded with TensorFlow -- the clean condition of Section~\ref{sec:lightingresults}.
The difference is \textbf{0.04 points}, the same size as the duplicate-row correction and well inside
the $\pm$0.31 seed spread. Pillow is not Android's decoder, so this measures what swapping one standard
libjpeg configuration for another does to accuracy, not the phone itself; it is nonetheless enough to
move the decoder from an anecdotally bounded risk to a measured, negligible one.

\textbf{The framing defect.} Every image in the corpus of Section~\ref{sec:construction} is a tight crop
around a single bunch. The application as first specified drew advisory framing brackets over roughly a
quarter of a full-bleed preview and classified the \textit{whole} frame, so a user obeying the on-screen
hint would hand the model a small, background-heavy bunch at a scale and context nothing in training
resembled. The parity fixtures cannot reveal this, because they are already crops. Re-scoring the full
2{,}722-image test split through simulated capture framings quantifies it (Table~\ref{tab:framing}).

\begin{table}[t]
\centering
\caption{Test-split accuracy under simulated capture framings. Full test split ($n{=}2{,}722$), seed-42
float16 TFLite artifact, Pillow decoding, synthetic backgrounds. Area fraction is the fraction of the
frame the bunch occupies, as achieved rather than as requested. Source:
\texttt{results/framing\_shift/} and \texttt{results/framing\_shift\_brackets/}.}
\label{tab:framing}
\begin{tabular}{@{}lrrrr@{}}
\toprule
Capture framing & Aspect & Area fraction & Accuracy (\%) & $\Delta$ (pts) \\
\midrule
Unmodified test crops (baseline) & --- & 1.00 & 94.16 & --- \\
\midrule
\textit{As specified:} whole frame classified & 9:20 & 0.25 & \textbf{48.46} & $-$45.70 \\
Squarer viewfinder, same area & 1:1 & 0.25 & 46.33 & $-$47.83 \\
Aspect distortion alone & 9:20 & 0.64 & 81.41 & $-$12.75 \\
Aspect distortion alone & 1:1 & 0.93 & 93.57 & $-$0.59 \\
\midrule
\textit{As built:} cropped to the brackets & 7:9 & 0.84 & \textbf{92.47} & $\mathbf{-1.69}$ \\
\bottomrule
\end{tabular}
\end{table}

At the framing as specified the model scores \textbf{48.46\%}, a \textbf{45.70-point} loss against the
same artifact's 94.16\% on unmodified crops. That is larger than the training-schedule correction of
Section~\ref{sec:schedule} and larger than the cross-source gap of Section~\ref{sec:lodo}: the largest
single effect in this report, and one produced entirely by user-interface geometry. It is also not
evenly distributed -- unripe and underripe recall collapse to 0.390 and 0.312 from 0.938 and 0.904, while
ripe and overripe degrade comparatively gently -- so the loss concentrates on the boundary this project
has found difficult everywhere else.

\textbf{Scale, not shape, is the binding term, which rules out the intuitive remedy.} Holding the bunch at
full frame and varying only the aspect ratio costs between 0.59 and 12.75 points; narrowing the bunch to a
quarter of the frame costs upward of 45 regardless of shape, and a square viewfinder at that area
(46.33\%) is marginally \textit{worse} than the tall one. Simply making the viewfinder squarer would
therefore not have addressed the defect.

\textbf{The remedy adopted, and what it recovers.} Of the available options -- leave the specification
alone, widen the brackets to most of the frame and rewrite the hint, or crop the capture to the brackets
before preprocessing -- the third was implemented. It was preferred over widening because it does not
depend on the user complying with an instruction: whatever is framed, the application cuts to the box the
user was aiming through, so the model's input is determined by code rather than by photographic
discipline. The brackets are measured against the camera preview at capture time and the decoded frame is
cropped to them, with the review screen displaying that same cropped region, so what the user approves is
what is classified. Scored at the shipped brackets' 7:9 aspect, the corrected framing reaches
\textbf{92.47\%} -- recovering \textbf{44.0 of the 45.7 points}. The residual 1.69 points falls almost
entirely on underripe (recall 0.904 to 0.834), with the other three classes within a point of baseline;
mean confidence in the predicted class also falls, from 0.805 to 0.761, which matters because the
application's close-call notice keys off that margin. Images supplied through the gallery are deliberately
not cropped: the application did not frame them and has no basis for locating the subject within them.

\textbf{What this does not establish.} The framing conditions are simulated. Backgrounds are synthetic --
a blurred continuation of the crop, or flat grey as a control -- rather than photographs of plantation
clutter, so these figures isolate the scale and aspect change faithfully while approximating everything
else. The correction's geometry is covered by unit tests, but \textbf{no photograph taken through the
corrected application has been scored, because field images could not be obtained}, and the recovery
figure therefore comes from the same simulator that found the defect rather than from an independent
observation. Two risks follow and are carried into Section~\ref{sec:limitations}: real background clutter
may not behave as the synthetic fill does, and a bunch overflowing the brackets is now clipped by the crop
where previously it was merely diluted by the surrounding frame. Finally, nothing here is a latency
measurement -- the instrumented run was an x86\_64 emulator, which says nothing about a phone.

\subsection{Lighting robustness}
\label{sec:lightingresults}

\begin{table}[t]
\centering
\caption{FastViT-T8 (seed 42) accuracy across the lighting stress-test conditions, for the trained Keras
checkpoint and for the converted float16 artifact the application ships. Full test split, 2,722 images per
condition.}
\label{tab:lighting}
\begin{tabular}{@{}lrr@{}}
\toprule
Condition & Keras checkpoint & Deployed float16 TFLite \\
\midrule
Clean test baseline & 94.56\% & 94.12\% \\
Shadow & 93.90\% & 93.57\% \\
Darkened ($\times$0.7) & 93.72\% & 93.79\% \\
Brightened ($\times$1.3) & 93.06\% & 92.76\% \\
\midrule
\textbf{Spread, this seed} & \textbf{1.51 pts} & \textbf{1.36 pts} \\
Spread, three-seed campaign & 1.02 $\pm$ 0.43 pts & -- \\
\bottomrule
\end{tabular}
\end{table}

\textbf{NFR-4 is met, and it holds for the artifact the application ships, not only for its Keras parent}
(Table~\ref{tab:lighting}). On the trained model the three-seed spread is 1.02 $\pm$ 0.43 points and no
condition costs more than 1.5 points (Figure~\ref{fig:lighting}). Re-run through the converted float16
model, the spread is 1.36 points against the same seed's 1.51: conversion does not make the model less
stable. That was not a foregone conclusion -- float16 halves every weight's mantissa, and darkening
compresses activations into the bottom of the value range where relative precision loss is largest -- yet
darkening is the condition that moves \textit{least} under conversion (+0.07 points). Every condition
clears the 85\% bar by more than seven points; underripe remains the weakest class (recall 0.877--0.903
across conditions) and no class collapses in any of them. The one-seed TFLite spread should not be compared
with the three-seed 1.02, which pools seed noise the single run does not carry. The ranking
across architectures is notable in its own right: FastViT-T8 1.02\,pts, EfficientNetB0 2.45\,pts,
MobileNetV2 5.19\,pts -- the selected architecture is not only the most accurate but by a clear margin the
most stable, and MobileNetV2's five-point spread is the widest measured. Brightening is consistently the
harshest condition, consistent with overexposure destroying the colour saturation cues on which the task
depends.

\begin{figure}[t]
\centering
\includegraphics[width=0.75\linewidth]{05_lighting_stress_test.png}
\caption{Accuracy across the four lighting conditions, by architecture.}
\label{fig:lighting}
\end{figure}

\subsection{Secondary experiments and negative results}
\label{sec:secondary}

\textbf{Ordinal loss (CORAL) -- negative result, mechanism identified.} Hypothesis: since ripeness is
ordinal, a rank-consistent ordinal head \citep{cao2020coral} should reduce underripe$\leftrightarrow$ripe
confusion better than categorical cross-entropy. Same backbone, same recipe, same class weighting; only
the loss and head changed. Result: 66.90\% against the categorical baseline's 73.73\% under the same (v1)
conditions. It \textit{did} reduce underripe$\rightarrow$ripe confusion from 34\% to 21\% and raised
underripe recall from 0.39 to 0.47, directionally confirming the hypothesis -- but it overcorrected:
ripe$\rightarrow$underripe confusion nearly quadrupled from 10\% to 34\%, ripe recall collapsed from 0.82
to 0.53, and lighting stability worsened from a 2.2-point spread to 6.1. The technique traded one
direction of the confusion for the other rather than resolving it, and was not adopted. This experiment
was run under recipe v1 and was not re-run under v2; the comparison is internally consistent but its
absolute numbers are superseded.

\textbf{Per-class threshold calibration -- a v1-era remedy the corrected recipe made redundant.} Per-class
decision thresholds were fitted on the validation split alone and applied unchanged to test
(Table~\ref{tab:calibration}). For the selected model \textbf{the gain (+0.21) is smaller than its own
seed-to-seed spread ($\pm$0.31)} and cannot be distinguished from run-to-run variance. The mechanism
explains why: calibration works by trading recall between adjacent classes, so its value depends on there
being a badly served class to trade toward. Under v1, FastViT's underripe recall was 0.510 and calibration
raised it to 0.681 at a cost of 0.120 in ripe recall; under v2 underripe arrives at 0.904 untouched, and
the same procedure moves it by 0.021. Roughly two-thirds of the fitted gain also fails to transfer (+0.60
validation against +0.21 test). \textbf{Decision: the headline stays uncalibrated at 94.55\%.} This
negative result corroborates Section~\ref{sec:schedule} from an independent direction: a technique
introduced specifically to compensate for poor underripe recall becomes superfluous once the training
schedule is fixed -- evidence that under-training, rather than a modelling or data limitation, was the
binding constraint.

\begin{table}[t]
\centering
\caption{Per-class threshold calibration gain, v1 vs.\ v2. Source: \texttt{calibration\_v2\_summary.md}.}
\label{tab:calibration}
\begin{tabular}{@{}lrrr@{}}
\toprule
Architecture & Test gain, v1 & Test gain, v2 & Validation gain, v2 \\
\midrule
EfficientNetB0 & +1.41 & +0.55 & +0.84 \\
\textbf{FastViT-T8 (selected)} & +0.45 & \textbf{+0.21} & +0.60 \\
MobileNetV2 & $-$0.24 & +0.29 & +0.46 \\
\bottomrule
\end{tabular}
\end{table}

\subsection{Measurement integrity: the duplicate-row defect}

The merged manifest's test split contains 2{,}722 rows over 2{,}695 unique images: 27 rows (0.99\%) are
exact duplicates originating in the bunch-piles crop pipeline. They never cross a split boundary, so this
is double-counting rather than leakage. The effect was measured directly rather than merely bounded
(Table~\ref{tab:duplicates}): inference was run once per unique image and every metric computed twice over
identical predictions -- once weighted by manifest multiplicity (reproducing what was published) and once
counting each image one.

\begin{table}[t]
\centering
\caption{Impact of the duplicate-row defect on reported accuracy. Source: \texttt{duplicate\_row\_impact.md}.}
\label{tab:duplicates}
\begin{tabular}{@{}lrrr@{}}
\toprule
Quantity & As published & Deduplicated & $\Delta$ \\
\midrule
FastViT-T8 test accuracy ($n{=}3$) & \textbf{94.55\% $\pm$ 0.31} & 94.51\% $\pm$ 0.33 & $-$0.04 pts \\
EfficientNetB0 (seed 42) & 91.55\% & 91.47\% & $-$0.08 pts \\
MobileNetV2 (seed 42) & 89.64\% & 89.61\% & $-$0.03 pts \\
\bottomrule
\end{tabular}
\end{table}

The largest shift across any run measured is \textbf{0.08 points} -- an order of magnitude below the
$\pm$0.31 seed spread the headline already carries, and far below the 9.55-point margin over NFR-1.
\textbf{No reported conclusion changes.} The bias direction is knowable rather than a coin flip: every
delta is negative because the 27 duplicated images are ones the models classify well (92--100\% on that
subset against 90--95\% overall), so double-counting inflates rather than deflates. What is measured is
the \textit{reporting} error, not the \textit{training} error -- the training split carries 1{,}422
duplicated rows whose effect would require retraining on a deduplicated manifest to quantify, which was
out of scope.

\section{Discussion and Conclusion}
\label{sec:discussion}

\subsection{Revisiting the research questions}

\begin{table}[t]
\centering
\caption{Research questions revisited.}
\label{tab:rq}
\begin{tabular}{@{}p{1.0cm}p{5.6cm}p{8.4cm}@{}}
\toprule
 & Question & Answer \\
\midrule
RQ1 & $\geq$85\% on a leakage-free multi-source split? & \textbf{Yes} -- all three architectures clear it; FastViT-T8 at 94.55\% $\pm$ 0.31\% \\
RQ2 & Which architecture, and is the margin real? & \textbf{FastViT-T8}, and yes -- but the margin over EfficientNetB0 is 1.82 points at matched depth, not the 3.29 implied by a naive comparison \\
RQ3 & Does it transfer to an unseen source? & \textbf{No} -- 52.11\% LODO mean, a 42.44-point gap, with a convention-fitting failure mode \\
RQ4 & Does it deploy within budget? & \textbf{Size yes} (6.1\,MB/100\,MB), \textbf{lighting yes} (1.02 pts spread), \textbf{application built} and faithful to the Python pipeline within decoder tolerance, \textbf{on-device latency undemonstrated} \\
\bottomrule
\end{tabular}
\end{table}

H1 and H3 are supported. H2 is supported at roughly half its apparent magnitude. H4 is half-resolved. The
project's most consequential finding was hypothesised by nobody: the binding constraint on accuracy was
the training schedule (Section~\ref{sec:schedule}), and correcting it moved every architecture by 15--16
points on unchanged data.

\subsection{What the generalization gap means}

The 42.44-point gap is the single most important number in this report for anyone considering field
deployment, and it deserves interpretation rather than apology. The main-split figure is measured on
unseen \textit{photographs} from sources the model also trained on. The LODO figure is measured on an
entirely unseen \textit{source} -- a different camera, a different site, and, decisively, a different
labelling convention. The gap is the part of main-split accuracy that does not survive that change.
Because ripeness is a continuous biological process cut into four categories by human convention
(Section~\ref{sec:related}), two datasets can label the same physical bunch differently without either
being wrong, and a model trained across several such datasets learns to reconcile them in whichever way
the training majority agrees on. Section~\ref{sec:lodo} shows exactly that: each held-out fold snaps to the
convention shared by the remaining three sources, and the class most sensitive to convention -- underripe,
the boundary case -- is the one that collapses.

Two things follow. First, \textbf{the gap is not primarily a capacity or a data-volume problem}, so it
would not obviously be closed by a larger model or by more images from the same four sources; it is a
label-semantics problem, and the appropriate remedies are label harmonisation against a single standard
(MPOB's, for a Malaysian deployment), fine-tuning on target-site data before use, or both. Second, \textbf{
it explains why the corrected recipe widened rather than narrowed the gap}: a properly trained model fits
its training sources' conventions harder. Under-training was inadvertently acting as regularisation
against convention-fitting -- an observation worth stating because it is counterintuitive and is directly
evidenced here.

The project's reporting decision follows from this analysis rather than from convenience. The NFR was
written against a held-out test split, which 94.55\% clears legitimately; the honest estimate of
performance at an unseen mill is 52.11\%; and both are stated wherever the headline appears. Stating the
limit unprompted is what makes the headline credible under examination.

\subsection{Limitations}
\label{sec:limitations}

\begin{itemize}
  \item \textbf{L1 -- The intended ViT could not be evaluated.} MobileViT was named in the plan and proved
  unloadable under Keras 3 across five integration routes; a sweep of fifteen lightweight transformer
  variants found only FastViT buildable. The limitation is a property of the current Keras 3 ecosystem,
  not of this project's configuration, and the substitution is documented rather than quiet.
  \item \textbf{L2 -- Fine-tuning depth is not exactly matched, and cannot be} (Section~\ref{sec:depth}).
  The residual 2.3-point mismatch favours FastViT, so 1.82 points is an upper bound on its architectural
  advantage. A $K{=}5$ control would bracket it and was not run.
  \item \textbf{L3 -- On-device latency is undemonstrated} (Section~\ref{sec:deployment}). Converted
  latency is three orders of magnitude inside the 2\,s budget, but on a development x86 CPU. No compliance
  claim is made on that basis.
  \item \textbf{L4 -- Cross-source generalization is weak and the LODO folds carry confounds}
  (Section~\ref{sec:lodo}): unequal fold training-set sizes, and a v1$\rightarrow$v2 comparison in which
  architecture and recipe changed together.
  \item \textbf{L5 -- The corpus is geographically skewed.} 69\% of rows come from one Indonesian
  mill-collected source; the only Malaysian source contributes 2.3\% and contains no underripe class at
  all. For a Malaysian deployment target this is a material limitation of the training distribution, and
  the Malaysia fold's 60.90\% -- the \textit{best} LODO fold -- is the closest thing this project has to a
  field estimate.
  \item \textbf{L6 -- Measured defects retained by decision, not oversight.} The duplicate-row defect moves
  the headline by 0.04 points and is reported rather than fixed; re-running the nine-run campaign to
  recover 0.04 points against a $\pm$0.31 spread was not justified. The fix belongs in the pipeline before
  any future campaign.
  \item \textbf{L7 -- CORAL and the FYP1 baseline were not re-run under recipe v2}, so their absolute
  numbers are superseded even though their internal comparisons hold.
  \item \textbf{L8 -- Numerical reproducibility.} The XLA autotuner emitted cuDNN algorithm-mismatch
  warnings during FastViT training and not for the CNNs. The three-seed protocol measures total
  run-to-run variance, which subsumes such nondeterminism, so the reported spreads already account for it.
  \item \textbf{L9 -- The deployed application cannot decode a JPEG the way training did, and the effect
  is measured as negligible} (Section~\ref{sec:ondevice}). Every accuracy figure in this report was
  measured on TensorFlow-decoded images; Android decodes through the platform's own libjpeg configuration,
  which disagrees by up to 16 levels of 255 per pixel and shifts individual probabilities by up to 0.042,
  with no predicted class changing on the twenty fixtures run on the emulator. Over the full 2,722-image
  test split, swapping TensorFlow's decoder for Pillow's moves the deployed artifact's accuracy by 0.04
  points (94.12\% $\rightarrow$ 94.16\%). Pillow stands in for Android's decoder rather than being it, so
  the claim is that a change of standard decoder does not measurably move accuracy; Android's own decoder
  has been exercised on the twenty fixtures only.
  \item \textbf{L10 -- The framing correction is verified in simulation, not in the field}
  (Section~\ref{sec:ondevice}). The 45.70-point framing defect and the 44.0 points the crop recovers are
  both measured by re-scoring the test split through simulated capture geometries with synthetic
  backgrounds. The correction's geometry is unit-tested, but no photograph taken through the corrected
  application has been scored, because field images could not be obtained within the project period. Two
  consequences remain unmeasured: whether real plantation background behaves as the synthetic fill does,
  and the new failure mode in which a bunch overflowing the brackets is clipped rather than diluted. This
  is the limitation most likely to separate this report's figures from field behaviour, and it is stated
  as an open risk rather than a closed one.
\end{itemize}

\subsection{Conclusion}

This project set out to determine whether a lightweight vision model could grade oil palm FFB ripeness
accurately enough, and cheaply enough, to run offline on a smartphone. On the terms the requirement was
written in, it can: FastViT-T8 reaches 94.55\% $\pm$ 0.31\% on a leakage-free four-source test split,
converts to a 6.1\,MB TensorFlow Lite artifact needing no non-standard runtime, and varies by only 1.02
points across daylight, shadow and low-light conditions. Four of the five non-functional requirements are
met; the fifth, on-device latency, is undemonstrated rather than failed, with converted inference
measuring 47\,ms against a 2{,}000\,ms budget on development hardware.

That is not the whole finding, and the report would be dishonest if it stopped there. The same model
scores 52.11\% when tested on a source it has never seen -- a 42.44-point gap that is the honest estimate
of performance at an unfamiliar mill. The analysis in Section~\ref{sec:discussion} locates the cause not
in model capacity but in the disagreement between datasets about where one ripeness category ends and the
next begins, which is a labelling problem that more training on the same corpus will not solve.

Three methodological results are, in the author's assessment, of wider value than the accuracy figures.
The first is that this project spent most of its duration reporting accuracies that were 15 points too low
because every run was stopped before convergence, and the defect was invisible in the results -- visible
only in the saved per-epoch histories, which is why they were saved. The second is that a controlled
ablation run on that broken pipeline returned a clean, plausible, and false result, and reversed when
re-run correctly. An ablation inherits the validity of the recipe it is run under. The third emerged only
once the model was put behind a user interface: a framing decision in the camera screen, invisible to
every model-side metric in this report and to the on-device parity tests alike, cost more accuracy than
any modelling choice examined here (Section~\ref{sec:ondevice}). A deployment path is not verified by
checking that the model was ported faithfully; it is verified by checking what the application actually
feeds it. All three were caught only because every dependent conclusion was re-derived rather than carried
forward, and all three are reported here as findings rather than as confessions.

\subsection{Future work}
\label{sec:futurework}

\begin{enumerate}
  \item \textbf{Complete the deployment path} -- benchmark the float16 artifact on a physical Android
  device against NFR-2. The application exists (Section~\ref{sec:ondevice}); what it has never run on is
  real hardware, which also leaves the torch control and the low-light thresholds untested. On ARM,
  integer quantization may outperform float16 via XNNPACK or NNAPI, which could reverse the float16
  recommendation.
  \item \textbf{Photograph bunches through the application and score them}, which is the one measurement
  that would convert the framing correction of Section~\ref{sec:ondevice} from a simulated result into an
  observed one, and would simultaneously test the two risks that simulation cannot reach: real plantation
  background, and bunches that overflow the framing brackets.
  \item \textbf{Attack the generalization gap at its cause.} Re-label a stratified subset of all four
  sources against a single standard (MPOB's) and re-measure LODO; if the gap narrows materially, the
  convention-mismatch diagnosis is confirmed and label harmonisation becomes the project's most valuable
  next investment.
  \item \textbf{Site fine-tuning.} Evaluate how many target-site images are needed to recover main-split
  accuracy at a new mill -- a directly actionable number for anyone deploying this.
  \item \textbf{The $K{=}5$ depth control} to bracket FastViT's fine-tuning depth from above and convert
  the 1.82-point upper bound into a point estimate.
  \item \textbf{Segmentation preprocessing}, cost-gated in two stages: first crop the existing test set
  with an off-the-shelf segmenter and re-evaluate the trained classifier -- if there is no lift, report a
  non-transferring result and stop; only on a real lift annotate 300--500 masks and train a YOLOv8n-Seg
  stage. Any accuracy gain must then be charged against the $\leq$2\,s and $\leq$100\,MB budgets.
  \item \textbf{Nine-angle crop augmentation} per \citet{suharjito2021mobile}, which this project did not
  adopt -- with source-photo split discipline preserved, since derived crops crossing splits is precisely
  the FYP1 leakage bug.
  \item \textbf{Re-run the FYP1 MobileNetV2 baseline and the CORAL experiment under recipe v2} for a
  like-for-like comparison, both being cheap.
\end{enumerate}

\section*{Acknowledgments}
The author thanks Dr.\ Muhammad Shahreeza Safiruz Kassim (supervisor) and Ahmad Syazwan Ramli, Malaysian
Palm Oil Board (co-supervisor), for their guidance throughout this project.

\bibliographystyle{unsrt}
\bibliography{references}

\appendix

\section{Figure index}

All figures are regenerated on the recipe-v2 campaign and live under \texttt{results/figures/} in the
project repository (Table~\ref{tab:figindex}).

\begin{table}[h]
\centering
\caption{Figure index and where each belongs in the report.}
\label{tab:figindex}
\begin{tabular}{@{}p{4.6cm}p{6.8cm}p{2cm}@{}}
\toprule
File & Content & Belongs in \\
\midrule
\texttt{01\_class\_distribution.png} & Per-class counts across the merged corpus & \S\ref{sec:construction} \\
\texttt{02\_source\_composition.png} & Per-source composition of each class & \S\ref{sec:construction} \\
\texttt{03\_confusion\_matrices.png} & Confusion matrices per architecture & \S\ref{sec:archcomparison} \\
\texttt{04\_per\_class\_metrics.png} & Per-class precision / recall / F1 & \S\ref{sec:archcomparison} \\
\texttt{05\_lighting\_stress\_test.png} & Accuracy across the four lighting conditions & \S\ref{sec:results} \\
\texttt{06\_overall\_summary.png} & Headline metric summary & \S\ref{sec:results} \\
\texttt{07\_training\_curves.png} & Per-epoch training and validation curves & \S\ref{sec:schedule} \\
\texttt{08\_underripe\_ripe\_confusion\_by\_source.png} & The hard boundary, broken down by source & \S\ref{sec:lodo} \\
\texttt{09\_multiseed\_architecture\_comparison.png} & Three-architecture comparison with seed spread & \S\ref{sec:archcomparison} \\
\texttt{10\_calibration\_effect.png} & Calibration gain, v2 refit & \S\ref{sec:secondary} \\
\texttt{11\_lodo\_generalization\_gap.png} & Main split vs.\ LODO, per fold & \S\ref{sec:lodo} \\
\texttt{framing\_shift/framing\_shift.png} & Accuracy against camera framing, by aspect and area fraction & \S\ref{sec:ondevice} \\
\bottomrule
\end{tabular}
\end{table}

\section{Run index}

Table~\ref{tab:runindex} maps each headline result to the run directory or aggregate artifact that
produced it, for reproducibility.

\begin{table}[h]
\centering
\caption{Reproducibility: result to source-artifact mapping.}
\label{tab:runindex}
\begin{tabular}{@{}p{4.6cm}p{8.8cm}@{}}
\toprule
Result & Run directory / artifact \\
\midrule
Headline, FastViT-T8 & \texttt{fastvit\_fyp2\_blockft\_v2\_gpu\_f50\_k3\_seed\{42,43,44\}/} \\
Architecture comparison & \texttt{multiseed\_v2\_summary.md} \\
Depth control & \texttt{effnet\_depth\_k4\_v2\_summary.md} \\
LODO campaign & \texttt{lodo\_v2\_summary.md} \\
TFLite export & \texttt{tflite\_export\_summary.md} \\
Calibration refit & \texttt{calibration\_v2\_summary.md} \\
Duplicate-row measurement & \texttt{duplicate\_row\_impact.md} \\
Framing defect (as specified) & \texttt{results/framing\_shift/framing\_shift.json} \\
Framing correction (as built) & \texttt{results/framing\_shift\_brackets/framing\_shift.json} \\
On-device parity & \texttt{app/src/androidTest/.../ModelParityTest.kt}, \texttt{reference.json} \\
Decisions log (full reasoning) & \texttt{PROJECT\_PLAN.md}, \S8 \\
\bottomrule
\end{tabular}
\end{table}

\end{document}
